% !TEX root = ../termpaper.tex
% @author Leonhard Lüdeke
%

\section{Architektur}

\subsection{Anforderungen für Architektur}

Die Zielarchitektur wird in dieser Arbeit nicht primär aus verfügbaren Frameworks abgeleitet, sondern aus fachlichen, organisatorischen und sicherheitsbezogenen Anforderungen entwickelt.

Daher ist der Ausgangspunkt nicht die Frage, was für Technologien verfügbar sind sondern welche Aufgaben das System erfüllen soll (User-Stories) und unter welchen Randbedingungen es betrieben wird. Hier sind sowohl Unternehmerischen Richtlinien (Richtlinie 27) wie auch Technische Begebenheiten zu berücksichtigen, sodass sich das Multi-Agenten-System nahtlos in die vorhandene Infrastruktur einfügt.

Durch die Analytische Betrachtung durch geeignete Threat-Modeling Prozesse werden weitere Designentscheidungen getroffen, welche fortlaufende Architekturentscheidungen begründen sollen.

Dies könnten zum Beispiel zusätzliche Kontrollmechanismen, Protokollentscheidungen oder auch stärkere zentralisierung/verteiltes Design des Multi-Agenten-Systems sein.


%\subsection{Threat Modeling}
%
%Mittels einer frühzeitiger Sicherheitsanalysen, werden besonders Schutzbedürftige Komponenten und Vertrauensgrenzen ausgemacht. \cite{shostack2014threat} Daraufhin können technische oder organisatorische Gegenmaßnahmen identifiziert werden.
%
%Da gerade in Agentischen Systemen Sicherheitsprobleme nicht isoliert an einzelnen Komponenten entstehen, sondern sich Risiken entlang von Kommunikationsketten, Delegationsentscheidungen, Tool-Interaktionen und Protokollbeziehungen entstehe. \cite{wong_adding_2000} Bietet sich eine breite Analyse mittels Threat-Modeling an um Sicherheitsrisiken auszumachen. \cite{owasp_owasp_nodate}

\subsection{Bisherige Überlegungen zur Architektur}

Die bisherige Zielarchitektur orientiert sich an einem modularen Multi-Agenten-Ansatz mit getrennten Verantwortlichkeiten. Ein Zentraler Orchestrator und/oder Workflow Controller erscheint für Multi-Agenten-Systeme eine plausible Kernkomponente. \cite{georgeff_communication_1983,jennings_roadmap_1998, vashishtha_cosmic_2026}

Dieser nimmt Anfragen entgegen, strukturiert Aufgaben, delegiert Teilaufgaben an spezialisierte Agenten und konsolidiert Ergebnisse. \cite{sycara_distributed_1996, xue_all--one_2024} Damit werden Koordination, Kontrolllogik und Fehlerbehandlung in dieser Instanz gebündelt, damit der Gesamtprozess transparent strukturiert werden kann. \cite{wong_adding_2000}

Als User-Interface ist eine Angepasste Version von OpenWebUI \cite{noauthor_open-webuiopen-webui_2026} vorgegeben an die sich das Multi-Agenten-System anbinden soll. Bisherige Agenten sind mittels Langchain und Langgraph \cite{chase_langchain_2022} in Python \cite{noauthor_welcome_2026} implementiert, sowie teilweise mit einem FastApi Gateway \cite{ramirez_fastapi_2026} mittels eines OpenApi Endpunktes \cite{noauthor_openapi_nodate} im internen Kubernetes cluster, \cite{noauthor_produktionsreife_nodate} in einem Docker container \cite{noauthor_docker_nodate} deployed.

Die bisherigen Überlegungen zur Architektonischen Umsetzung, sowie bisher maßgeblich beteiligte Systeme sind im angegebenen Diagramm dargestellt. \ref{fig:mas_arch}
